\subsection{Bayesian Normal Comparisons}
\label{subsec:bayesian-normal-comparisons}
% Source integration: attached Chapter 11 screenshots, Section 11.2.4.

The normal two-sample model also has a Bayesian interpretation.  With diffuse priors, the resulting posterior formulas resemble frequentist $t$ intervals, but the conditioning and interpretation are different.

\begin{definition} \label{def:bayesian-common-variance-two-sample-model} \hlt{Common-Precision Two-Sample Normal Model}
\begin{enumerate}[label=(\roman*)]
\setlength{\itemsep}{0pt}
\item \hlt{Common-Precision Model}: $\theta=(\mu_X,\mu_Y,\xi)\in\R^2\times(0,\infty)$ with shared precision $\xi=\sigma^{-2}$ and variance $\xi^{-1}$.
\item \hlt{Conditional model}: Given $(\mu_X,\mu_Y,\xi)$, the variables $X_1,\ldots,X_n,Y_1,\ldots,Y_m$ are independent, with $X_i\mid(\mu_X,\xi)\sim N(\mu_X,\xi^{-1})$ and $Y_j\mid(\mu_Y,\xi)\sim N(\mu_Y,\xi^{-1})$.
\item \hlt{Improper prior}: $\pi(\mu_X,\mu_Y,\xi)\propto \xi^{-1}$ on $\mu_X,\mu_Y\in\R$ and $\xi>0$.
\item \hlt{Observed data}: Write $x=(x_1,\ldots,x_n)$ and $y=(y_1,\ldots,y_m)$.
\item \hlt{Sample means}: $\overline{x}\defeq n^{-1}\sum_{i=1}^{n}x_i$ and $\overline{y}\defeq m^{-1}\sum_{j=1}^{m}y_j$.
\item \hlt{Sample variances}: $s_X^2\defeq(n-1)^{-1}\sum_{i=1}^{n}(x_i-\overline{x})^2$ and $s_Y^2\defeq(m-1)^{-1}\sum_{j=1}^{m}(y_j-\overline{y})^2$.
\item \hlt{Pooled variance}: $s_p^2\defeq\{(n-1)s_X^2+(m-1)s_Y^2\}/(m+n-2)$.
\item \hlt{Posterior kernel}: $\pi(\mu_X,\mu_Y,\xi\mid x,y)\propto \xi^{(m+n)/2-1}\exp\{-(\xi/2)[\sum_{i=1}^{n}(x_i-\mu_X)^2+\sum_{j=1}^{m}(y_j-\mu_Y)^2]\}$.
\end{enumerate}
\end{definition}

\begin{proposition} \label{prop:bayesian-two-sample-conditional-posterior} \hlt{Conditional Posterior Factors}\\
For $n,m\geq2$, put $\nu\defeq m+n-2$ and take $\operatorname{Gamma}$ shape-rate.\\
In the common-precision model, $\pi(\mu_X,\mu_Y,\xi\mid x,y)=\pi(\mu_X\mid\xi,x,y)\pi(\mu_Y\mid\xi,x,y)\pi(\xi\mid x,y)$,\\
where $(\mu_X,\mu_Y)\mid\xi,x,y\sim N(\overline{x},(n\xi)^{-1})\otimes N(\overline{y},(m\xi)^{-1})$ and $\xi\mid x,y\sim\operatorname{Gamma}(\nu/2,\nu s_p^2/2)$.\\
Hence $\mu_X\perp\!\!\!\perp\mu_Y\mid(\xi,x,y)$.
\begin{proof} Let $S_X(\mu_X)\defeq\sum_{i=1}^{n}(x_i-\mu_X)^2$, $S_Y(\mu_Y)\defeq\sum_{j=1}^{m}(y_j-\mu_Y)^2$, and $A\defeq(n-1)s_X^2+(m-1)s_Y^2=\nu s_p^2$.\\
Completing squares gives $S_X(\mu_X)=(n-1)s_X^2+n(\mu_X-\overline{x})^2$ and $S_Y(\mu_Y)=(m-1)s_Y^2+m(\mu_Y-\overline{y})^2$.\\
Thus $\pi(\mu_X,\mu_Y,\xi\mid x,y)\propto\xi^{(m+n)/2-1}\exp\{-\xi A/2\}\exp\{-n\xi(\mu_X-\overline{x})^2/2\}\exp\{-m\xi(\mu_Y-\overline{y})^2/2\}$.\\
For fixed $\xi$, last two factors are independent normal kernels for $\mu_X$, $\mu_Y$ with variances $(n\xi)^{-1}$, $(m\xi)^{-1}$.\\
Also $\int\exp\{-n\xi(\mu_X-\overline{x})^2/2\}\,d\mu_X\propto(n\xi)^{-1/2}$ and $\int\exp\{-m\xi(\mu_Y-\overline{y})^2/2\}\,d\mu_Y\propto(m\xi)^{-1/2}$, so marginalising $\mu_X,\mu_Y$ leaves $\pi(\xi\mid x,y)\propto\xi^{\nu/2-1}\exp\{-\xi\nu s_p^2/2\}$.\\
\hlt{Structural Commentary}: Completing squares isolates location kernels and pooled residual $A=\nu s_p^2$.\\
\hlt{Key Insight}: Marginalising $\mu_X,\mu_Y$ leaves the pooled gamma law for $\xi$.
\end{proof}
\end{proposition}

\begin{proposition} \label{prop:bayesian-two-sample-t-posterior} \hlt{Posterior $t$ Law for a Mean Difference}\\
In the common-precision model with $n,m\geq2$, set $\Delta\defeq\mu_X-\mu_Y$, $d\defeq\overline{x}-\overline{y}$, $c\defeq n^{-1}+m^{-1}$, $\nu\defeq m+n-2$, and $T\defeq(\Delta-d)/(s_p\sqrt{c})$. Then $T\mid x,y\sim t_\nu$.
\begin{proof} By Proposition \ref{prop:bayesian-two-sample-conditional-posterior}, $\Delta\mid\xi,x,y\sim N(d,c/\xi)$ and $\xi\mid x,y\sim\operatorname{Gamma}(\nu/2,\nu s_p^2/2)$.\\
Equivalently, take independent $E\sim N(0,1)$, $V\sim\chi^2_\nu$, $\xi=V/(\nu s_p^2)$, and $\Delta=d+\sqrt{c/\xi}\,E$.\\
Thus $T=(\Delta-d)/(s_p\sqrt{c})=E/\sqrt{V/\nu}$, so $T\mid x,y\sim t_\nu$.\\
\hlt{Structural Commentary}: Normal given $\xi$; Student after mixing over $\xi$.\\
\hlt{Key Insight}: The $t$ law is exactly a standard normal divided by an independent estimated-scale factor $\sqrt{V/\nu}$.
\end{proof}
\end{proposition}

\begin{remark} \label{rem:bayesian-two-sample-interpretation} \hlt{Same Algebra, Different Conditioning}\\
The statistic has the same algebra as the pooled two-sample $t$ pivot in Theorem \ref{thm:two-sample-pooled-t-pivot}: frequentist inference treats $\Delta$ as fixed and $X,Y$ as random, whereas the posterior treats $x,y$ as fixed and $\Delta\mid x,y$ as random.
\end{remark}

\begin{example} \label{ex:bayesian-heat-fusion-comparison} \hlt{Latent Heat Posterior Comparison}\\
For method contrast $\Delta=\mu_X-\mu_Y$, the common-precision posterior gives $\mathbb{P}(\Delta>0\mid X,Y)\approx0.998$ and posterior interval $\mathbb{P}(0.015\leq\Delta\leq0.065\mid X,Y)\approx0.95$, numerically matching the pooled normal-theory interval.
\end{example}
